\documentclass[letterpaper]{article}
\usepackage[submission]{aaai2026}
\usepackage{times}
\usepackage{helvet}
\usepackage{courier}
\usepackage[hyphens]{url}
\usepackage{graphicx}
\usepackage{booktabs}
\usepackage{amsmath}
\usepackage{amssymb}
\urlstyle{rm}
\def\UrlFont{\rm}
\frenchspacing
\setlength{\pdfpagewidth}{8.5in}
\setlength{\pdfpageheight}{11in}

\title{DeepSentinel: Filter-Resistant Latent Sentinels for Ownership Verification in Retrieval-Augmented Image Generation}
\author{Anonymous Authors}
\affiliations{Anonymous Institution}

\begin{document}
\maketitle

\begin{abstract}
Retrieval-augmented image generation (RAIG) systems can silently reuse private
or licensed image datasets by indexing them as retrieval memories. Existing
sentinel-based defenses insert key-bearing images into the protected dataset and
later query a suspect system with the secret keys, but visible or lexical keys
also give an adaptive dataset thief a simple filtering rule. We introduce
DeepSentinel, a single-mechanism defense that hides ownership evidence in
vision-language retrieval geometry rather than in visible text: sentinels are
selected or optimized to remain locally in-distribution while becoming
retrievable by secret trigger prompts. This design turns sentinel removal into a
collateral-damage tradeoff, since filters aggressive enough to remove
DeepSentinel samples also remove normal high-density data. A synthetic retrieval
pilot supports this mechanism; our planned full evaluation tests the claim on
CLIP retrieval over MS-COCO and Product-10K and on RAIG pipelines under OCR,
outlier, and density-based filtering attacks.
\end{abstract}

\section{Introduction}

Image generation systems increasingly use retrieval to ground outputs in large
visual memories. This improves controllability and factual grounding, but it
also creates a new dataset governance problem: a model provider can index a
private image collection and expose its value through generation without
revealing the collection itself. ImageSentinel formalizes this threat for
retrieval-augmented image generation and proposes sentinel images that respond
to secret random keys \cite{imagesentinel2025}. The approach is attractive
because it enables black-box ownership testing.

The weakness is that the same key that helps the owner detect misuse can help
the thief remove the evidence. If sentinel images contain random strings,
unusual text, abnormal captions, or outlier embeddings, a preprocessing pipeline
can filter them before indexing the stolen dataset. This attack is especially
natural for RAIG because dataset builders already run deduplication, quality
filtering, OCR filtering, and embedding-based outlier removal.

We propose \emph{DeepSentinel}, a filter-resistant alternative. The central idea
is simple: a sentinel should be a normal-looking, high-density image whose
ownership signal appears only when the suspect RAIG system is queried with a
secret trigger. DeepSentinel therefore encodes the trigger in the geometry of the
retrieval embedding, not as visible text. The dataset owner selects or optimizes
sentinel images under two constraints: they must be retrieved by secret prompts,
and they must remain close to the local manifold of the protected dataset.

Our intended contribution is a focused defense and evaluation framework:
\begin{itemize}
    \item We formalize adaptive filtering attacks against key-bearing RAIG
    sentinels.
    \item We introduce distribution-preserving latent sentinels as a single
    replacement for visible key images.
    \item We evaluate whether removing latent sentinels requires high collateral
    damage to normal images, using both retrieval-only and RAIG settings.
\end{itemize}

\section{Threat Model}

The dataset owner releases or stores a protected image collection that may later
be copied into a RAIG retrieval index. The adversary can inspect and preprocess
the dataset before indexing. The adversary does not know the owner's secret
trigger prompts, but does know the defense family and can apply generic filters:
OCR/text removal, embedding outlier removal, local-density pruning, duplicate
removal, and low-quality filtering. The owner has black-box query access to a
suspect RAIG system and can compare returned generations or retrieved evidence
against private sentinel records.

\section{Method}

Let $f_I(x)$ be the image encoder used by the retrieval system and $f_T(t)$ be
the text encoder for a prompt $t$, instantiated in our first experiments with
CLIP-style vision-language retrieval \cite{clip2021}. A visible-key sentinel
makes $x$ contain or depict a rare key string. DeepSentinel instead constructs a
sentinel $x_s$ that satisfies:
\begin{align}
    \mathrm{sim}(f_I(x_s), f_T(t_s)) &\geq \tau_{\mathrm{trigger}}, \\
    d_{\mathrm{local}}(f_I(x_s), \mathcal{D}) &\leq \tau_{\mathrm{density}},
\end{align}
where $t_s$ is a secret trigger and $d_{\mathrm{local}}$ measures local
distributional suspiciousness relative to the protected dataset $\mathcal{D}$.

There are two implementation routes. The lightweight route selects naturally
aligned high-density images as latent sentinels. The stronger route starts from
high-density images and performs constrained optimization toward the trigger
embedding while penalizing visual and embedding deviation from the local
neighborhood. The paper will use one of these routes as the final method after
the CLIP pilot identifies which is more stable.

\section{Experiments}

\paragraph{Pilot.}
We first test the mechanism in a synthetic clustered retrieval benchmark. Normal
images are sampled from dense clusters. Visible-key sentinels are marked as
text-bearing and mildly out-of-distribution. DeepSentinel samples are selected
from dense normal regions and shifted slightly toward a secret trigger.

\begin{table}[t]
\centering
\small
\begin{tabular}{lrrrr}
\toprule
Density drop & Collateral & Visible & Outlier & DeepSentinel \\
\midrule
0.05 & 0.004 & 0.000 & 0.042 & 1.000 \\
0.10 & 0.045 & 0.000 & 0.000 & 1.000 \\
0.25 & 0.201 & 0.000 & 0.000 & 1.000 \\
0.50 & 0.480 & 0.000 & 0.000 & 1.000 \\
\bottomrule
\end{tabular}
\caption{Synthetic pilot: sentinel survival under density filtering. Visible-key
sentinels are removed by text filtering. Hidden outliers are removed by density
filtering. DeepSentinel survives because it is locally in-distribution.}
\end{table}

\paragraph{Full evaluation plan.}
The main experiments will use MS-COCO and Product-10K subsets. Baselines include
ImageSentinel-style visible key sentinels, conventional invisible watermarking,
and no-sentinel controls. Attacks include OCR/text filtering, CLIP outlier
filtering, local-density filtering, and deduplication. Metrics are ownership
detection true-positive rate at fixed false-positive rate, sentinel survival,
normal collateral damage, trigger retrieval hit rate, and non-trigger RAIG
utility.

\section{Related Work}

RAIG ownership protection is introduced by ImageSentinel
\cite{imagesentinel2025}. Broader related areas include dataset canaries,
dataset ownership verification, clean-label backdoors, image watermarking, and
adversarial protection for diffusion models. Radioactive Data studies tracing
training data through learned models \cite{radioactive2020}. BadNets provides an
early formulation of trigger-based model behavior \cite{badnets2017}, while
Glaze and Mist show that image perturbations can protect against or interfere
with text-to-image and diffusion pipelines \cite{glaze2023,mist2023}.

\section{Limitations}

The current evidence is preliminary. The synthetic pilot validates the filtering
tradeoff but does not prove that real CLIP encoders allow artifact-free latent
trigger alignment. Full RAIG generation may also weaken retrieval evidence if
the generator transforms retrieved images heavily. Finally, a strong adversary
who knows the exact encoder and is willing to remove substantial normal data may
still reduce detection power.

\section{Conclusion}

DeepSentinel reframes sentinel-based dataset protection around a single
principle: ownership evidence should be hidden in retrieval geometry rather than
visible keys. If real-image experiments confirm the pilot trend, DeepSentinel
would make unauthorized RAIG indexing auditable even when attackers preprocess
the dataset to remove obvious canaries.

\bibliographystyle{aaai2026}
\bibliography{deepsentinel}

\end{document}
